\documentclass[11pt]{article}
\usepackage[a4paper,margin=27mm]{geometry}
\usepackage{amsmath,amssymb,amsthm,mathtools}
\usepackage{graphicx}
\usepackage[hidelinks]{hyperref}
\usepackage{xcolor}

\newtheorem{theorem}{Theorem}
\newtheorem{lemma}{Lemma}
\newtheorem{corollary}{Corollary}
\newcommand{\A}{\mathcal A}
\newcommand{\Bnet}{\mathcal B}
\newcommand{\F}{\mathcal F}
\newcommand{\UCP}{\operatorname{UCP}}
\newcommand{\QuOp}{\operatorname{QuOp}}
\newcommand{\Prob}{\operatorname{P}}
\newcommand{\Ad}{\operatorname{Ad}}
\newcommand{\Extr}{\operatorname{Extr}}

\title{Poincar\'e-Covariant Quantum Operations\\as a Double-Coset Hypergroup}
\author{Candidate solution packet for arXiv:2303.15171, Problem 5.6}
\date{27 August 2026}

\begin{document}
\maketitle

\begin{abstract}
Giorgetti asks whether the relative quantum operations of a Haag-dual full
Poincar\'e-covariant subnet $\Bnet\subset\A$ form the double-coset hypergroup
$G/\!/H$ supplied by Doplicher--Roberts reconstruction and subsystem
classification.  We prove the assertion in precisely that standard regime.
The local discrete-subfactor theorem identifies the maps on each double cone
with the probability measures on $G/\!/H$.  Compatibility makes the local
representing measures equal, and gauge--Poincar\'e commutation gives the
converse construction.  Thus
\[
  \UCP(\A\mid\Bnet)\cong \Prob(G/\!/H),\qquad
  \QuOp(\A\mid\Bnet)\cong G/\!/H
\]
affinely, topologically, and with the compact-hypergroup operations.
\end{abstract}

\section{The source problem and its intended scope}

The source is Luca Giorgetti, \emph{Quantum Operations in Algebraic QFT},
arXiv:2303.15171~\cite{Gio}.  Problem 5.6 on printed page 9 reads as follows.

\begin{center}
 \includegraphics[width=0.98\textwidth]{figures/open_problem_crop.png}
\end{center}

In transcription: if $(\A,U,\Omega)$ is a Poincar\'e-covariant net of local
observables and $\Bnet\subset\A$ is a Haag-dual full Poincar\'e-covariant
subnet, prove
\[
        \QuOp(\A\mid\Bnet)\cong G/\!/H,
\]
where $G$ is the Doplicher--Roberts compact gauge group of $\Bnet$ and
$H<G$ is closed.

The notation in the problem already invokes Doplicher--Roberts reconstruction
and the Carpi--Conti subsystem classification.  Those theorems have standing
hypotheses beyond the short list in the survey's Definition 5.1.  We therefore
state the exact structural conclusion that the proof uses.  This both covers
the intended classified setting and prevents those extra hypotheses from
being silently assumed.  In particular, the result applies whenever the
classification cited immediately before Problem 5.6 supplies
\begin{equation}\label{eq:fixedpoint}
     \Bnet=\F^G\subset \F^H=\A,
\end{equation}
with its standard local irreducibility, discreteness, and locality properties
for the type $\mathrm{III}$ inclusions.  Carpi--Conti obtain precisely this
fixed-point description under their standing AQFT assumptions
\cite[Theorem 5.2]{CC}; compare the discussion in
\cite[Example 9.26]{BDG}.

\section{Idea of the proof}

For a fixed double cone $O$, the hard operator-algebraic work is already
contained in the local theorem of Bischoff--Del Vecchio--Giorgetti: all UCP
maps on $\F(O)^H$ fixing $\F(O)^G$ are uniquely parametrized by probability
measures on $G/\!/H$~\cite[Section 9.2, Proposition 9.7]{BDG}.

The apparent global obstruction is that a coherent family could conceivably
choose a different measure $\mu_O$ for every $O$.  If $O_1\subset O_2$,
however, compatibility identifies the restrictions of the two maps on
$\A(O_1)$.  The local parametrization at $O_1$ is injective, so
$\mu_{O_1}=\mu_{O_2}$.  Any two double cones are contained in a third; hence
one measure parametrizes the entire family.  Conversely, averaging the global
gauge automorphisms against one fixed bi-$H$-invariant measure is automatically
compatible and vacuum preserving.  Internal gauge symmetries commute with
$U$, so this family is Poincar\'e covariant.  Taking extreme points changes
the probability simplex into its Dirac masses, namely $G/\!/H$ itself.

\section{The local parametrization}

Let $K:=G/\!/H=H\backslash G/H$, let $q:G\to K$ be the quotient map, and let
$m_H$ be normalized Haar measure on $H$.  A probability measure $\mu$ on
$K$ has the canonical bi-$H$-invariant lift
\begin{equation}\label{eq:lift}
 \widetilde\mu
  :=\int_K m_H*\delta_g*m_H\,d\mu(HgH)\in\Prob(G).
\end{equation}
The right-hand side is independent of the representative $g$.

Write $\alpha:G\to\operatorname{Aut}(\F)$ for the gauge action.  For a double
cone $O$ and $x\in\A(O)=\F(O)^H$, set
\begin{equation}\label{eq:theta}
 \Theta_O(\mu)(x):=\int_G\alpha_g(x)\,d\widetilde\mu(g).
\end{equation}
Bi-$H$-invariance makes the output $H$-fixed, while every element of
$\Bnet(O)=\F(O)^G$ is fixed pointwise.  Proposition 9.7 of~\cite{BDG}, applied
to
\[
       \F(O)^G\subset\F(O)^H,
\]
states that
\begin{equation}\label{eq:localhomeo}
 \Theta_O:\Prob(K)\longrightarrow
 \{\psi:\A(O)\to\A(O):\psi\text{ is normal UCP and }
                       \psi|_{\Bnet(O)}=\mathrm{id}\}
\end{equation}
is an affine homeomorphism.  In this local discrete setting the maps on the
right are automatically faithful, vacuum preserving, and vacuum-adjointable;
see the UCP characterization preceding the compact-hypergroup construction
in~\cite[Sections 6--7]{BDG}.  The same proposition identifies composition
and vacuum adjoint with the double-coset convolution and involution.

Because the gauge action is a net action, for $O_1\subset O_2$ one has the
crucial naturality identity
\begin{equation}\label{eq:naturality}
 \Theta_{O_2}(\mu)|_{\A(O_1)}=\Theta_{O_1}(\mu).
\end{equation}

\section{The global theorem}

Equip $\UCP(\A\mid\Bnet)$ with the pointwise ultraweak topology used in the
source: its members are compatible, vacuum-preserving, Poincar\'e-covariant
families of normal faithful UCP maps that fix $\Bnet$ pointwise.

\begin{theorem}\label{thm:main}
Assume the fixed-point realization \eqref{eq:fixedpoint}, where $G$ acts
faithfully and minimally by unbroken internal symmetries of the complete field
net $\F$, $H<G$ is closed, and each inclusion
$\F(O)^G\subset\F(O)^H$ has the irreducible local discrete type
$\mathrm{III}$ properties used in Proposition 9.7 of~\cite{BDG}.  Then
\begin{equation}\label{eq:globalhomeo}
 \Theta:\Prob(G/\!/H)\longrightarrow\UCP(\A\mid\Bnet),
 \qquad \Theta(\mu):=\{\Theta_O(\mu)\}_{O},
\end{equation}
is an affine homeomorphism.  It intertwines convolution and involution.
Consequently
\[
       \QuOp(\A\mid\Bnet)\cong G/\!/H
\]
as a compact hypergroup.
\end{theorem}

\begin{proof}
We prove existence, surjectivity, topology, and extremality in turn.

\smallskip\noindent
\emph{The integral construction gives a global family.}
Equation~\eqref{eq:theta} is normal, unital, and completely positive.  Gauge
automorphisms preserve the vacuum state, hence so does their average.  The
local vacuum state is faithful, so a vacuum-preserving positive map cannot
annihilate a nonzero positive element; thus the map is faithful.  It fixes
$\Bnet(O)$ pointwise.  Naturality~\eqref{eq:naturality} gives compatibility.

The action of the unbroken internal gauge group commutes with the Poincar\'e
representation:
\[
 \Ad U(p)\circ\alpha_g=\alpha_g\circ\Ad U(p)
 \qquad(p\in\mathcal P,\ g\in G).
\]
Using this identity under the integral gives
\[
 \Ad U(p)\circ\Theta_O(\mu)\circ\Ad U(p)^{-1}
      =\Theta_{pO}(\mu),
\]
so $\Theta(\mu)\in\UCP(\A\mid\Bnet)$.

\smallskip\noindent
\emph{Every global family has one representing measure.}
Let $\Phi=\{\phi_O\}_O\in\UCP(\A\mid\Bnet)$.  By the local affine
homeomorphism~\eqref{eq:localhomeo}, for each $O$ there is a unique
$\mu_O\in\Prob(K)$ such that
\[
                    \phi_O=\Theta_O(\mu_O).
\]
If $O_1\subset O_2$, compatibility and~\eqref{eq:naturality} give
\[
 \Theta_{O_1}(\mu_{O_1})=\phi_{O_1}
 =\phi_{O_2}|_{\A(O_1)}
 =\Theta_{O_1}(\mu_{O_2}).
\]
Injectivity of $\Theta_{O_1}$ yields $\mu_{O_1}=\mu_{O_2}$.  Given arbitrary
double cones $O_1,O_2$, choose a double cone $O_3$ containing both.  The
previous argument gives $\mu_{O_1}=\mu_{O_3}=\mu_{O_2}$.  Thus all $\mu_O$
equal one measure $\mu$, and $\Phi=\Theta(\mu)$.  This proves surjectivity;
injectivity follows already from any single $\Theta_O$.

\smallskip\noindent
\emph{Topology and operations.}
For every $O$, the coordinate map $\Theta_O$ is continuous, so $\Theta$ is
continuous for the product point-ultraweak topology.  Fixing one double cone
$O_0$, the inverse is the composition of evaluation at $O_0$ with
$\Theta_{O_0}^{-1}$, hence is continuous.  Therefore~\eqref{eq:globalhomeo}
is an affine homeomorphism.  The local result gives
\[
 \Theta_O(\mu*\nu)=\Theta_O(\mu)\circ\Theta_O(\nu),
 \qquad
 \Theta_O(\mu^*)=\Theta_O(\mu)^{\sharp}.
\]
Equality in every double cone proves the corresponding global identities.

\smallskip\noindent
\emph{Extreme points.}
The extreme points of $\Prob(K)$ are exactly the Dirac measures, so the
extreme points of the relative convex set are $\Theta(\delta_k)$, $k\in K$.
It remains only to match this with the source convention
\[
 \QuOp(\A\mid\Bnet)
 =\QuOp(\A)\cap\{\Phi:\Phi|_{\Bnet}=\mathrm{id}\}.
\]
The relative set is a face of $\UCP(\A)$.  Indeed, suppose
$\Phi=t\Phi_1+(1-t)\Phi_2$, $0<t<1$, and $\Phi$ fixes $\Bnet$.  For every
unitary $u\in\Bnet(O)$,
\[
       u=t\phi_{1,O}(u)+(1-t)\phi_{2,O}(u).
\]
The two terms on the right are contractions, while a unitary is an extreme
point of the unit ball of a unital $C^*$-algebra.  Hence
$\phi_{1,O}(u)=\phi_{2,O}(u)=u$.  Since every element of a unital
$C^*$-algebra is a linear combination of unitaries, both $\Phi_i$ fix
$\Bnet$.  Thus extremes of the relative set are precisely its elements that
are extreme in $\UCP(\A)$, completing the identification with $K$.
\end{proof}

\begin{corollary}[Problem 5.6 in the classified AQFT regime]
Under the standing Doplicher--Roberts and Carpi--Conti hypotheses invoked in
Problem 5.6, its proposed conclusion holds:
\[
                   \QuOp(\A\mid\Bnet)\cong G/\!/H.
\]
In particular, it is compact and carries the asserted double-coset
hypergroup structure.
\end{corollary}

\begin{proof}
The cited subsystem classification supplies~\eqref{eq:fixedpoint}, with $G$
the Doplicher--Roberts gauge group of $\Bnet$ and $H<G$ closed
\cite[Theorem 5.2]{CC}.  The associated local fixed-point subfactors are the
double-coset examples of~\cite[Section 9.2 and Example 9.26]{BDG}.  Apply
Theorem~\ref{thm:main}.
\end{proof}

\section{Verification, limitations, and novelty check}

The proof has no numerical component.  Its only substantive external inputs
are (a) the fixed-point realization from subsystem classification and (b) the
local affine homeomorphism of Proposition 9.7 in~\cite{BDG}.  Once these are
in place, the new local-to-global step uses only compatibility, injectivity,
and the fact that double cones are directed.  The face argument above also
checks explicitly that taking relative extreme points agrees with the source
definition.

The scope is important.  If one reads only the minimal axioms in the survey's
definition of a local net, without the standard assumptions required for
Doplicher--Roberts reconstruction and the Carpi--Conti classification, the
fixed-point presentation need not be available and the group $G$ asserted in
the problem may not be defined.  The packet makes no claim for that broader
reading; it proves the problem in the very classified regime named in its
statement and preceding paragraph.

A bounded novelty search on 27 August 2026 covered the run's local indexes and
live exact-title, exact-phrase, arXiv-id, author/keyword, and
``$\QuOp(\A\mid\Bnet)$'' searches.  It found the source paper, the conformal-net
antecedent, and the local-subfactor theorem, but no later paper claiming the
Poincar\'e-covariant global result.  This is evidence for novelty, not an
exhaustive literature certification.

The recommended human check is the precise transfer from the standing
hypotheses intended by Problem 5.6 to the local hypotheses of Proposition 9.7.
The accompanying \texttt{verification.md} gives a line-by-line dependency
audit.

\begin{thebibliography}{9}

\bibitem{Gio}
L. Giorgetti,
\emph{Quantum Operations in Algebraic QFT},
arXiv:2303.15171 (2023), Problem 5.6.

\bibitem{BDG}
M. Bischoff, S. Del Vecchio, and L. Giorgetti,
\emph{Compact hypergroups from discrete subfactors},
J. Funct. Anal. \textbf{281} (2021), 109004;
arXiv:2007.12384.

\bibitem{CC}
S. Carpi and R. Conti,
\emph{Classification of subsystems for graded-local nets with trivial
superselection structure},
Comm. Math. Phys. \textbf{253} (2005), 423--449;
arXiv:math/0312033.

\end{thebibliography}

\end{document}
